\documentclass[10pt,a4paper]{amsart}
\usepackage[margin=19mm]{geometry}
\usepackage{amsmath,amssymb,mathtools}
\usepackage[scr=rsfs,cal=euler]{mathalfa}
\usepackage{xfrac}
\usepackage[unicode,hidelinks,bookmarks=false]{hyperref}
\newcommand{\ie}{i.e.\ }
\newcommand{\cf}{cf.\ }
\newcommand{\resp}{respectively\ }
\newcommand{\Subsection}{\subsection}
\providecommand{\nobreakdash}{\nobreak\hbox{-}}
\setlength{\emergencystretch}{2em}
\multlinegap0pt
\usepackage{amssymb}
\usepackage{algorithm}\usepackage{algpseudocode}
\newtheorem{proposition}{Proposition}
\newcommand{\real}{\mathbb{R}}
\newcommand\xbf{\mathbf{x}}
\title{\LARGE \bf Induced Stackelberg Equilibrium Seeking via Iterative Tikhonov Regularization}
\author{Silvia Cianchi, Anibal Sanjab, Sergio Grammatico}
\date{Source: arXiv:2604.26746v1 (CC BY 4.0)}
\begin{document}
\maketitle

\noindent\emph{Source and licence.} \LARGE \bf Induced Stackelberg Equilibrium Seeking via Iterative Tikhonov Regularization. Version arXiv:2604.26746v1 (CC BY 4.0), distributed by arXiv under the Creative Commons Attribution 4.0 International licence, http://creativecommons.org/licenses/by/4.0/, as recorded in the arXiv licence metadata for this exact version and shown on the abstract page https://arxiv.org/abs/2604.26746v1. Full licence notice: \texttt{licenses/arxiv-2604.26746-CC-BY-4.0.txt}.\par\smallskip

\noindent\textit{Editorial scope.} Counted unit: the article's proposition Cost/Osc (source0.tex lines 286--288). The article's complete original proof of it is delivered with it (source0.tex lines 289--291, one \texttt{proof} environment of 3 lines). No mathematics is written, repaired or re-derived: the statement and the proof are the article's own lines, in the article's own notation, and no retained line is substituted or re-flowed.  Retained uncounted as the source-local context the proof needs: lines 257--263; lines 270--272; lines 281--283.  Nothing was omitted from any retained span, so the package carries no omission record.  No retained line contains a citation, so the wrapper carries no bibliography and no bibliography entry.  No retained line contains a figure environment, an image-inclusion command or drawing code, so no image is delivered and the package declares no asset dependency.  Editorial formatting only, in addition: an AMS \texttt{amsart} preamble that restates the article's own declarations and macros used by the retained lines, namely the environments \texttt{proposition} on separate counters, together with the standard packages those lines need.  The retained environments print in this document as Proposition 1 in order of appearance, whereas the article prints the same items with the numbering of its own full document.  The complete native source of the article as distributed in the arXiv e-print for this exact version is bundled unchanged as \texttt{sources/199418/source0.tex}.
\begin{equation}\label{Stackelberg-ILL}
\begin{aligned}
\underset{y} {\min} \, & J_0(y,\xbf):=y^2+y(x_1+x_2)
\\
\mathrm{s. t.} \, \,  &\xbf\in \mathcal{E}^\mathrm{int}(y).
\end{aligned}
\end{equation}

\begin{align}\label{update}
    y_{k+1}=y_k-\eta \nabla_k,
\end{align}

\begin{align}\label{inexactNabla}
\nabla_k=2y_k(1-x_{1 \mid k})+x_{1 \mid k}(1-\epsilon).
\end{align}

\begin{proposition}\label{Cost/Osc}
    Let us assume that $\left(y_k\right)_{k \in \mathbb{N}}$ generated by \eqref{update} converges, i.e., there exists $y_{\infty} \in \real$ such that $\lim \limits_{k \rightarrow \infty} y_k = y_{\infty}$. Then, there exists $ x_{1 \mid \infty}\in \real, \, x_{1 \mid \infty} \neq 1$, such that $\lim \limits_{k \rightarrow \infty} x_{1 \mid k} = x_{1 \mid\infty}$. Furthermore, $y_\infty$ is the SE solution of the game in \eqref{Stackelberg-ILL} with $\xbf=\mathrm{col}(x_{1 \mid\infty}, \, -(y+\epsilon)x_{1 \mid \infty})$.
\end{proposition}

\begin{proof}
    Consider the update rule in \eqref{update} with $\nabla_k$ as in \eqref{inexactNabla}. If $y_k \rightarrow y_{\infty} \in \real$, then necessarily $y_{k+1}-y_k \rightarrow 0$, which implies $2 y_{k}(1-x_{1 \mid k})+x_{1 \mid k}(1-\epsilon) \rightarrow 0$. Taking the limiting value, it holds that $\frac{- x_{1|k}(1-\epsilon)}{2(1-x_{1 \mid k})} \rightarrow y_\infty$. For this limit to exist and be finite, it necessarily holds $\lim \limits_{k \rightarrow \infty}\frac{x_{1 \mid k}}{1-x_{1 \mid k}}\in \real$, which implies $x_{1 \mid k} \rightarrow x_{1 \mid \infty}\in \real,$ for some $ x_{1 \mid \infty} \neq 1$.
\end{proof}

\end{document}
